\begin{textsample}{NEG Dim 4 – ro – Score -2.00 – ro/ro\_em\_16.txt}  \label{ex:f4_neg_049}
Quadro 2.
Proposição de exemplos de estratégias de metodologias ativas para o trabalho pedagógico com o estudante no Ensino Médio em Rondônia | Tipo de Estratégia | O que é? | Quando usar? | Como usar? | Como avaliar? | |---|---|---|---|---| | Aprendizagem colaborativa | Consiste em uma estratégia ativa de aprendizagem que requer a participação, integração e colaboração de todos os membros da equipe. | Pode ser usada em processos de aprendizagem com vistas a estimular os estudantes a fazer e receber críticas sobre temáticas estudadas.
E ainda, para o trabalho em equipe, a tomada de decisão e a melhora na comunicação. | Em ambientes de sala de aula ( presencial e virtual ).
Criar grupos e atribuir funções aos integrantes conforme as possibilidades e o interesse coletivo.
Exemplos : contrato de aprendizagem ; diagrama dos cinco porquês ; giro colaborativo ; jogos de cartas ; passa ou repassa ; disputa argumentativa ; aprendizagem em espiral. | Envolver os grupos de modo que todos possam apontar os pontos positivos e de atenção da participação de cada grupo. | | Aprendizagem entre pares ou times ( TBL ) | Caracteriza -se como uma estratégia de aprendizagem entre pares e times na qual os estudantes aprendem e ensinam ao mesmo tempo. | Em situações educativas voltadas para a troca e compartilhamento de informações e/ou a construção colaborativa e participativa de ideias sobre a temática estudada. | Em ambientes de sala de aula ( presencial e virtual ).
Criar grupos e atribuir funções aos integrantes conforme as possibilidades e o interesse coletivo.
Exemplos : estudo individual/coletivo, entrevista, conferência, filmes, experimentos, realização de tarefas ( individual/coletivo ) ; apelação, feedback ; estudos de caso ; aplicação de testes ( individual/coletivo ). | Envolver os grupos de modo que todos possam apontar os pontos positivos e de atenção da participação de cada grupo. | | Exposições interativas | Trata -se de uma estratégia ativa de docência interativa que ocorre mediante relações dialógicas e participativas entre aluno e professor. | É utilizada quando o professor deseja problematizar e dinamizar a aula, e deste modo, envolver o estudante em situações de aprendizagem concretas sobre a temática abordada. | A exposição interativa não ocorre somente em salas de aula equipadas com recursos tecnológicos, mas também pode ser planejada para sala de aula não informatizada.
Exemplos : debate, intercâmbio com o autor ; mapa mental, mural de fatos e notícias ; paleta de cores ; uso de aplicativos. | A partir da interação e participação do estudante nas dinâmicas desenvolvidas, como, por exemplo, em atividades como jogos de perguntas e respostas ; músicas ; filmes. | | Aprendizagem baseada em projetos ou problemas | Consiste em uma estratégia de aprendizagem que alia a construção do conhecimento de forma ativa com vista à solução colaborativa de desafios chamados problemas. | Recomenda -se utilizar quando se busca que o estudante explore soluções dentro de um contexto específico.
E ainda, é utilizado para despertar o lado inventivo, crítico e colaborativo do estudante ao participar de projetos diversificados. | A estratégia pode ser desenvolvida em espaços de aprendizagem diversificados.
Entretanto, é essencial o trabalho em grupos de forma colaborativa.
Exemplos : análise de fatores ou ideias ; estudo de caso ; uso de aplicativos ; construção de situações-problema ; matriz de problemas ; mapeamento de causas ; árvore dos problemas ; problemas do cotidiano. | Envolver os grupos de modo que todos possam apontar os pontos positivos e de atenção da participação de cada grupo.
Em paralelo, avaliar a capacidade de tomada de decisão, definição das ações, os prazos e o trabalho em grupo. | Fonte : adaptado de Costa ( 2020 ) ; Camargo ( 2018b ). | Tipo de Estratégia | O que é? | Quando usar? | Como usar? | Como avaliar? | |---|---|---|---|---| | Design Thinking | É uma estratégia de ensino-aprendizagem ativa, criativa e colaborativa.
Também é conhecida como uma metodologia para solução de problemas e/ou uma abordagem que catalisa a colaboração, a inovação e a busca por soluções. | É utilizada em situações de aprendizagem em que se deseja : {\EmojiFont ▪} Propor processos educativos para conversas ; iniciar, aprender com os erros, falhas e resolver discordâncias ; {\EmojiFont ▪} Definir um desafio estratégico ( problema ) ; {\EmojiFont ▪} Avaliar conhecimentos preexistentes sobre a realidade do estudante ; {\EmojiFont ▪} Planejar em conjunto soluções para um problema específico ; {\EmojiFont ▪} Identificar capacidades e perfis para implementar soluções ; {\EmojiFont ▪} Desenvolver habilidades do estudante referente ao gerenciamento da projeção de uma solução para um problema complexo ; {\EmojiFont ▪} Transformar ideias em realidade. | O Design Thinking pode ser desenvolvido em espaços de aprendizagem diversificados, mediante a observação e a cocriação, a partir do conceito de prototipagem ( criação e testes de protótipos ) rápida e análise de diferentes realidades. {\EmojiFont ▪} Pode ser utilizado diversos materiais ( papelão, massa de modelar ou papel ). {\EmojiFont ▪} São exemplos : storyboards, diagramas, mockup, filmagens, peças de teatro, maquetes, páginas da internet ; brainstorm com post-its. | {\EmojiFont ▪} Observe os conhecimentos prévios dos estudantes sobre os problemas levantados ; {\EmojiFont ▪} Analise os protótipos ( material ) produzidos pelos grupos.
O foco no processo de aprendizagem é constatar o que e como o estudante está aprendendo ; {\EmojiFont ▪} Oportunize momentos de diálogos para \textbf{verificação} se o aluno aprendeu o que estava previsto nos objetivos de aprendizagem. | | Gamificação | {\EmojiFont ▪} É uma técnica que usa elementos, design e pensamento de jogos em diferentes contextos para gerar engajamento, mudanças de comportamento, interatividade, resolução de problemas, atividades mais divertidas, motivar a ação e promover a aprendizagem do estudante. | {\EmojiFont ▪} Recomenda -se o uso quando se busca : a resolução de problemas, principalmente sociais ; o desenvolvimento de habilidades complexas como : tomada de decisão, desenvolvimento de liderança e inovação ; mudança de hábitos ( exemplo : inserção da atividade física na rotina ; aprender um idioma ). {\EmojiFont ▪} Pode ser usado também quando se pretende alcançar objetivos de aprendizagem mensuráveis e definidos por regras de forma interativa e com a presença de feedback. | {\EmojiFont ▪} Pode ser desenvolvida em espaços de aprendizagem diversificados que contemplem elementos como interatividade, ações individuais, coletivas, a estruturação de regras e o feedback. {\EmojiFont ▪} É importante no game : conhecer os objetivos/metas de aprendizagem ; definir os comportamentos e tarefas que fazem parte da solução ; conhecer seus jogadores ; selecionar o tipo de conhecimento ensinado ; assegurar a presença da diversão ; e utilizar as ferramentas apropriadas ; fazer protótipos e testagem. {\EmojiFont ▪} Exemplos : jogos pedagógicos ; jogos eletrônicos e populares ; jogos de aplicativos. | {\EmojiFont ▪} Elabore mapas de empatia para conhecer o tipo de jogador ; {\EmojiFont ▪} Verifique os avanços de aprendizagem a partir da \textbf{taxonomia} de Bloom ; {\EmojiFont ▪} Verifique o grau de dificuldade de cada game proposto ; {\EmojiFont ▪} Observe a motivação, a ação e ofereça feedback constante da atividade gamificada ; {\EmojiFont ▪} Utilize os pontos, níveis, fases e placar para dar feedback do game. | Fonte : adaptado de Cavalcanti e Filatro ( 2016 ) ; Flora ( 2015 ) ; Camargo ( 2018b ).

% matched lemmas: taxonomia, verificação
\end{textsample}
